\documentclass[10pt,a4paper]{article}
\usepackage[left=2.2cm,right=2.2cm,top=2.0cm,bottom=2.0cm]{geometry}
\usepackage{fontspec}
\usepackage{amsmath,amssymb}
\usepackage{booktabs}
\usepackage{graphicx}
\usepackage{tabularx}
\usepackage{array}
\usepackage{caption}
\usepackage{enumitem}
\usepackage[hidelinks]{hyperref}
\newcolumntype{Y}{>{\raggedright\arraybackslash}X}
\captionsetup{font=small,labelfont=bf,labelsep=space}
\setlength{\parindent}{2em}
\setlength{\parskip}{0pt}
\setlength{\tabcolsep}{3pt}
\emergencystretch=3em
\title{Prediction--Planning Consistent LSTM-Enhanced MPC-DWA for UAV Dynamic Obstacle Avoidance}
\author{}
\date{}
\begin{document}
\maketitle
\begin{abstract}
\noindent\textbf{Abstract. }

Dynamic obstacles can invalidate local UAV planning when observation, prediction and control are temporally misaligned. We introduce an LSTM trajectory predictor into a coupled model predictive control (MPC) and dynamic window approach (DWA) framework. The predictor uses the latest 1 s obstacle history to forecast the next 3 s, and the same time-indexed forecast is supplied to MPC safety costs and DWA clearance scoring. Training uses 800 independent trajectories, a 640/160 trajectory-level split, training-only normalization, aligned noise/delay augmentation, early stopping and best-validation selection. LSTM succeeds in 12/12 fixed trials and 36/36 random trials in the clean simulation suite, with fixed-window ADE/FDE of 0.49/1.14 m. Under clean, noise-only, delay-only and joint pressure tests, success rates are 2/3, 2/3, 2/3 and 1/3. Per-run outcomes, failure types and timing decomposition are reported, together with a classical Kalman baseline. The results support improved prediction–planning consistency under the tested conditions, but not universal robustness to arbitrary sensor delay or real-flight deployment.

\textbf{Keywords:} UAV; dynamic obstacle avoidance; model predictive control; dynamic window approach; LSTM; Kalman filter; trajectory prediction

\end{abstract}

\section{Introduction}

Unmanned aerial vehicles operating among moving objects must simultaneously achieve goal arrival, dynamically feasible motion and collision avoidance. Static-obstacle planning can treat occupied geometry as fixed constraints, whereas the future positions of dynamic obstacles change while a local trajectory is being executed. A local planner that relies only on the latest observation can therefore become unsafe before the short-horizon control action is completed. Constant-velocity extrapolation is also systematically biased when an obstacle turns, accelerates, pauses or switches between waypoints.

Model predictive control (MPC) can optimize a finite-horizon trajectory under dynamics and input constraints, while the dynamic window approach (DWA) selects executable controls from a locally reachable control set. Combining the two allows a planner to use both predictive optimization and short-horizon feasibility. The benefit, however, depends on whether future obstacle positions are supplied to both layers with the same time index. If MPC anticipates a turn while DWA evaluates candidates against a frozen or constant-velocity obstacle, the two layers can generate internally inconsistent decisions.

This paper asks whether a shared learning-based obstacle predictor, synchronously injected into MPC and DWA, can improve closed-loop arrival and collision-related outcomes relative to static freezing, constant-velocity extrapolation and a Kalman baseline in fixed three-dimensional scenes and extended random scenes. The contributions are:

\begin{enumerate}[leftmargin=*]
\item A closed-loop LSTM-MPC-DWA architecture that predicts 30 future steps from a 10-step history and sends the time-aligned prediction to both the MPC safety cost and the DWA temporal clearance score.
\item A trajectory-level training and validation protocol based on 800 independent trajectories, together with a classical Kalman prediction baseline and three matched closed-loop random seeds.
\item An evaluation suite containing fixed scenes, 12 independently generated random mixed-maneuver scenes, per-run outcomes, failure types, timing decomposition, structural ablations and four sensing-pressure conditions.
\end{enumerate}

The conclusions are deliberately bounded. The evidence concerns the tested simulation dynamics, generated scenes and sensing conditions; it does not establish universal superiority under arbitrary prediction errors, sensor delays or real-flight conditions.

\section{Related Work}

Hierarchical MPC-DWA planning connects medium-horizon model prediction with short-horizon executable control and is suitable for systems with velocity, angular-rate and dynamic constraints [1–4]. Classical DWA uses a dynamic window and short-horizon trajectory scoring to maintain executable candidates [5, 6]. MPC generates a receding-horizon reference and can incorporate constraints and safety costs [11]. Reviews of learning-based MPC emphasize that data-driven models should be considered jointly with constraints and online optimization [12].

Dynamic-obstacle prediction can be based on static freezing, constant-velocity models, Kalman filtering or learned sequence models. The Kalman filter provides a classical recursive state-estimation baseline for linear dynamical systems [7]. Static freezing is inexpensive but does not represent temporal motion; constant-velocity extrapolation is useful for near-linear motion but cannot describe acceleration, turning, speed variation or waypoint switching. LSTM models can extract local motion patterns from temporal histories [8] and have been used for interactive trajectory prediction [9]. Prediction-enhanced DWA [4] and physics-prior LSTM-DWA methods [10] further suggest that a forecast becomes useful for closed-loop avoidance only when its time indexing is consistent with the planner.

The present work does not treat the predictor as a complete avoidance method in isolation. Instead, it studies a specific coupling mechanism in which one predicted obstacle sequence is shared by the MPC planning layer and the DWA execution layer. This design makes it possible to separate offline prediction accuracy from closed-loop planning behavior and to report cases in which the two rankings do not coincide.

\section{Method}

\subsection{UAV Dynamics and Closed-Loop Timing}

The simulation step is dt = 0.1 s. The UAV state and control are defined as

\[
\mathbf{x}=[x,y,z,\theta,\psi,v]^\mathsf{T},\quad
\mathbf{u}=[v_c,\omega_\theta,\omega_\psi]^\mathsf{T}.
\]

The UAV uses a first-order velocity-response model:

\[
\begin{aligned}
x_{t+1}&=x_t+dt\,v_t\cos\theta_t\cos\psi_t,\\
y_{t+1}&=y_t+dt\,v_t\cos\theta_t\sin\psi_t,\\
z_{t+1}&=z_t+dt\,v_t\sin\theta_t,\\
\theta_{t+1}&=\theta_t+dt\,\omega_{\theta,t},\\
\psi_{t+1}&=\psi_t+dt\,\omega_{\psi,t},\\
v_{t+1}&=v_t+dt(v_{c,t}-v_t)/\tau,
\end{aligned}
\]

where τ = 0.5 s. The initial state is [1, 1, 1]ᵀ and the goal is [18, 18, 8]ᵀ. The simulated arena is [0, 20] × [0, 20] × [0, 10] m and the maximum run length is 650 steps. The input limits are 0 ≤ v\_c ≤ 2.0 m s−1, |ω\_θ|, |ω\_ψ| ≤ 0.8 rad s−1, with maximum acceleration 1.0 m s−2, pitch/yaw acceleration limits 0.5 rad s−2 and maximum deceleration 1.0 m s−2. A run reaches the goal when the three-dimensional UAV–goal distance is below 0.5 m and the complete trajectory contains no collision.

\subsection{LSTM Obstacle-Trajectory Prediction and Training Protocol}

The same network is applied independently to each dynamic obstacle. The input is the latest T\_h = 10 steps of an eight-dimensional feature vector:

\[
[x,y,z,v_x,v_y,v_z,\cos\psi,\sin\psi].
\]

The network outputs three-dimensional displacement increments for the next N\_p = 30 steps, corresponding to a 3 s horizon. Absolute predicted positions are obtained by adding the increments to the latest position in the input history. The training set contains 800 independently generated trajectories covering straight motion, turning, helix, acceleration, hover-then-maneuver behavior, sinusoidal speed variation, zigzag motion and waypoint following. The trajectories yield 25,569 time-aligned input–target windows; after augmentation, 40,910 training sequences and 10,228 validation sequences are used.

The split is performed at the trajectory level, with 640 trajectories for training and 160 for validation. The fixed S1–S4 scenes and the random evaluation trajectories are excluded from training. Means and standard deviations for inputs and targets are estimated only from training windows and then fixed for validation and testing. The network has an eight-dimensional sequence input, a 96-unit LSTM layer that returns the last output, dropout of 0.2, a 96-unit fully connected ReLU layer and a 90-dimensional regression output.

Training uses Adam with a learning rate of 10−3, batch size 128, a maximum of 180 epochs, L2 regularization of 10−4 and a gradient threshold of 1.0. Each training window receives one additional history view with position perturbation of standard deviation 0.03 m and one one-step delayed-history view. For the noisy view, the displacement target is recomputed relative to the newest perturbed observation, preventing a label shift caused by absolute-position noise. The validation set retains clean windows and adds an identically distributed 0.03 m pressure view. Validation loss is evaluated every 50 updates; training stops after 12 validations without improvement and restores the best-validation-loss model. The training seed is 20260905, and the split and training metadata are stored with the model. The selected model metadata records OutputNetworkIteration = 5150, with final validation loss 3.4524295 and RMSE 2.6277099 in normalized target-space units; these are not metre-scale prediction errors.

Under noisy or delayed observation, the closed loop additionally uses a causal position-smoothing window and delay compensation. Smoothing uses only current and past samples and preserves the latest observed position. For a delayed history, the already-lagged first d prediction steps are discarded and the tail is filled using the terminal predicted velocity. A complete noisy observation stream is generated once at the beginning of each run and is shared by all MPC and DWA prediction calls.

\subsection{Bidirectional MPC-DWA with Shared Prediction}

The MPC generates a reference trajectory over an N = 30-step horizon. For obstacle j at future step k, let the predicted position be \(\hat{\mathbf{o}}_{j,k}\). The signed clearance is

\[
d_{j,k}=\|\mathbf{p}_k-\hat{\mathbf{o}}_{j,k}\|-r_j-r_e,
\]

where r\_j is the obstacle radius and r\_e = 0.3 m is the UAV radius. At planning time t, \(\mathbf{p}_{t+k|t}\) and \(\hat{\mathbf{o}}_{j,t+k|t}\) denote the planned UAV position and obstacle prediction at the same future index. In the implementation, prediction row k is aligned with planning state X(:, k+1).

The MPC obstacle cost is

\[
J_{obs}=\sum_{k,j}\left[w_{obs}\max(0,-d_{j,k})^2
+w_s\max(0,d_s-d_{j,k})^2\right],
\]

where w\_obs = 20,000, w\_s = 12,000 and the soft safety boundary is d\_s = 0.35 m. A three-dimensional cross-track soft cost to the start–goal line is added to prevent permanent lateral displacement after repeated detours; the main experiment uses w\_route = 3.

The complete MPC objective is

\[
J=J_{goal}+J_{ctrl}+J_{smooth}+J_{route}+J_{obs},
\]

\[
\begin{aligned}
J_{goal}&=w_{pos}\sum_{k=0}^{N-1}\|\mathbf p_k-\mathbf g\|_2^2+w_{term}\|\mathbf p_N-\mathbf g\|_2^2,\\
J_{ctrl}&=w_{ctrl}\sum_{k=0}^{N-1}\|\mathbf u_k\|_2^2,\\
J_{smooth}&=w_{smooth}\sum_{k=1}^{N-1}\|\mathbf u_k-\mathbf u_{k-1}\|_2^2,\\
J_{route}&=w_{route}\sum_{k=1}^{N}\left\|\mathbf e_k-\frac{\mathbf e_k^{\mathsf T}\mathbf a}{\|\mathbf a\|_2^2}\mathbf a\right\|_2^2,
\quad \mathbf e_k=\mathbf p_k-\mathbf s,\ \mathbf a=\mathbf g-\mathbf s.
\end{aligned}
\]

The implementation uses w\_pos = w\_term = 1, w\_ctrl = 10 and w\_smooth = 50; s and g are the start and goal positions.

DWA constructs a dynamic window from the current control, maximum acceleration and maximum angular rate. Candidate controls are sampled around the first MPC control while retaining an exploration component. Each candidate is checked for short-horizon braking feasibility and predicted minimum clearance. Candidates with predicted clearance not exceeding 0.25 m are rejected. If no candidate satisfies the safety test, the candidate with the smallest predicted overlap is selected within the dynamic window and a fallback step is recorded. MPC and DWA use the same time-indexed obstacle sequence. The executed DWA control is fed back as a safety control box for the next MPC solve, producing bidirectional MPC-DWA coupling.

\subsection{Implementation and Reproducibility Parameters}

All methods use the same UAV dynamics, obstacle trajectories, MPC cost, DWA dynamic window and solver settings. Only the future-obstacle interface differs. MPC uses CasADi with IPOPT, a maximum of 300 iterations and a tolerance of 10−6. The base DWA sampling number is 300 and can be increased adaptively according to prediction confidence. The guided-sampling ratio is 0.35 and the exploration intensity is 3. The sample count is N\_s = round[N\_base(1 + α exp(−βτ))], with N\_base = 300, α = 1 and β = 3. Guided samples are drawn within radius 0.35 of the first MPC control and the remaining samples explore the dynamic window uniformly. Trust uses d\_safe = 1 m, d\_sense = 5 m, α\_T = 3 and σ = (0.35, 0.30, 0.15, 0.20). The score weights use a sigmoid with slope K = 8 and midpoint τ\_0 = 0.5: w\_pos ∈ [0.10, 0.30], w\_head ∈ [0.08, 0.22], w\_vel = 0.12(1 + sigmoid), and w\_obs ∈ [0.20, 0.45], followed by normalization; an additional vertical-corridor weight of 0.25 discourages unnecessary vertical escape. The software environment is MATLAB R2025a, CasADi 3.8.0 and the Deep Learning Toolbox. Fixed-scene seeds are 1000 × run seed + scene index; random-scene seeds are 100000 × run seed + scene index; sensing-pressure seeds are 700000 + 10000 × condition index + run seed. The scene-generation and training seeds are both 20260905.

\begin{table}[t]
\centering
\small
\caption{Main implementation and reproducibility parameters}
\label{tab:table-1}
\begin{tabularx}{\linewidth}{YYY}
\toprule
Parameter & Setting & Role \\
\midrule
Simulation step dt & 0.1 s & Closed-loop update period \\
MPC/prediction horizon N, N\_p & 30 steps / 3 s & Common planning and prediction horizon \\
History length T\_h & 10 steps / 1 s & LSTM input window \\
LSTM & 96 units, dropout = 0.2 & Sequence encoder \\
Training & 800 trajectories; 640/160 trajectory split; Adam; 180 epochs maximum; batch = 128; learning rate = 10−3 & Training and validation \\
Augmentation and selection & 0.03 m history perturbation; one-step delay; L2 = 10−4; gradient threshold 1.0; early stopping after 12 validation checks & Robustness and overfitting control \\
DWA sampling & N\_base = 300; r\_guide = 0.35; α = 1, β = 3 & Candidate generation \\
Safety parameters & r\_e = 0.3 m; d\_s = 0.35 m; d\_reject = 0.25 m; z ∈ [0.5, 9.5] m; nominal-height band = 3.0 m & Clearance and vertical-escape control \\
Route-recovery soft cost & w\_route = 3 & Suppression of persistent lateral offset \\
Control limits & v\_c ∈ [0, 2.0] m s−1; angular-rate limit 0.8 rad s−1; acceleration limit 1.0 m s−2; angular-acceleration limit 0.5 rad s−2; maximum deceleration 1.0 m s−2 & Dynamic feasibility and braking \\
Environment and termination & Start [1, 1, 1] m; goal [18, 18, 8] m; arena 20 × 20 × 10 m; maximum 650 steps & Scene bounds and termination \\
Solver & CasADi + IPOPT; max\_iter = 300; tol = 10−6 & MPC numerical solution \\
\bottomrule
\end{tabularx}
\end{table}

\section{Experimental Design}

\subsection{Evaluation Scenes and Compared Methods}

S1–S4 are four fixed three-dimensional scenes. The number of obstacles increases from 2, 4 and 5 to 6. Obstacle radii are approximately 1.2–2.0 m, and all obstacles cross the nominal route from [1, 1, 1]ᵀ to [18, 18, 8]ᵀ. The motion patterns include straight motion, sinusoidal speed variation, acceleration, turning, zigzag motion, hover-then-maneuver behavior and waypoint following. Obstacle trajectories are fixed before each run and do not read or respond to the UAV state.

The compared methods are defined as follows. static freezes each obstacle at its latest observed position. cv extrapolates with the latest observed velocity. kalman uses a constant-velocity Kalman state model and outputs future positions. lstm uses the proposed LSTM prediction and injects it synchronously into MPC and DWA.

Each fixed scene is run with matched seeds 1–3, yielding 12 trials per method and 48 fixed-scene trials in total. A separate scene seed, 20260905, generates 12 random scenes, each containing 3–8 obstacles and mixed maneuver types. The random scenes use the same three matched run seeds for all four methods, yielding 144 additional closed-loop trials. Randomness comes from DWA candidate sampling; the scene, dynamics, network parameters and controller weights are held constant within each matched comparison.

\subsection{Metrics and Independent Units}

Offline prediction is evaluated with average displacement error (ADE) and final displacement error (FDE). The fixed-window evaluation contains 17 obstacles and 20 time windows per obstacle, giving 340 windows per method. These windows overlap in time and are therefore treated as descriptive prediction measurements rather than 340 independent replications.

Closed-loop metrics include success rate, failure type, collision time steps, collision events, true minimum clearance, path length, three-dimensional tracking RMSE relative to the nominal route, DWA fallback steps and computation time per step. For obstacle j at time t, e\_j,t,k = ||ô\_j,t+k|t − o\_j,t+k||\_2. The implementation averages the N\_p future-step errors for each obstacle and time, then averages over visible obstacles and closed-loop times; FDE uses the same aggregation at k = N\_p:

\[
ADE=\frac{1}{T}\sum_t\frac{1}{M_t}\sum_{j=1}^{M_t}\frac{1}{N_p}\sum_{k=1}^{N_p}e_{j,t,k},\qquad
FDE=\frac{1}{T}\sum_t\frac{1}{M_t}\sum_{j=1}^{M_t}e_{j,t,N_p}.
\]

True minimum clearance is d\_min = min\_t,j(||p\_t − o\_j,t||\_2 − r\_j − r\_e); negative values indicate geometric overlap. A collision step has d\_j,t < 0, and a collision event is a transition from non-collision to collision. Success is 1(||p\_end − g||\_2 < 0.5 m and collision steps = 0). Computation time is decomposed into prediction, DWA candidate processing, MPC solution and state-update/metric computation. Step-level P95 identifies long-tail delay and excludes the initial t = 0 MPC solve. The primary independent unit is the matched scene–seed–method closed-loop run; time steps within a run are not treated as independent observations for significance testing.

\subsection{Additional Experiments}

To diagnose the source of closed-loop gains, a single-seed structural ablation is conducted for w\_route = 0 and 3, LSTM injection into MPC only, LSTM injection into DWA only, injection into both layers, and H\_inj = 10, 20 and 30. This experiment is used for mechanism diagnosis and does not replace the multi-seed main comparison.

S4 is additionally evaluated under clean observation, 0.05 m position noise, one-step delay and joint noise-plus-delay conditions. Each method is run with three matched seeds under each condition. The complete noisy observation stream is generated once per run and reused by all prediction calls. The suite contains one fixed six-obstacle S4 scene, 4 conditions × 3 seeds × 4 methods = 48 run-level records. Because n = 3 per method and condition, the 95\% Wilson intervals are wide; for example, 2/3 corresponds to 20.8–93.9\% and the joint LSTM result of 1/3 corresponds to 6.1–79.2\%. These tests define the boundary under sensing non-idealities and are not mixed with the clean-observation main experiment or interpreted as a multi-scene robustness estimate.

\section{Results}

\subsection{Fixed-Window Prediction}

\begin{table}[t]
\centering
\small
\caption{Fixed history–future window prediction errors}
\label{tab:table-2}
\begin{tabularx}{\linewidth}{YYYY}
\toprule
Method & Windows & ADE (m) & FDE (m) \\
\midrule
static & 340 & 2.07 & 3.92 \\
cv & 340 & 0.51 & 1.23 \\
kalman & 340 & 0.56 & 1.31 \\
lstm & 340 & 0.49 & 1.14 \\
\bottomrule
\end{tabularx}
\end{table}

Across fixed windows, LSTM reduces ADE and FDE by approximately 3.9\% and 7.3\% relative to cv and by 12.5\% and 13.0\% relative to Kalman. The S1–S4 ADE/FDE values for LSTM are 0.576/1.203, 0.363/0.833, 0.256/0.603 and 0.751/1.760 m, respectively. The corresponding Kalman values are 0.433/0.958, 0.367/0.913, 0.478/1.067 and 0.803/1.903 m. The LSTM advantage is most apparent in the nonlinear corridor-compression scene S3; it is not lowest for every window in S1. This supports a conditional nonlinear-prediction benefit but does not independently establish closed-loop safety.

\subsection{Multi-Seed Closed-Loop Comparison}

\begin{table}[t]
\centering
\small
\caption{Fixed-scene closed-loop summary over three matched seeds (mean ± standard deviation)}
\label{tab:table-3}
\begin{tabularx}{\linewidth}{YYYYYYYY}
\toprule
Method & Runs & Success rate (95\% Wilson CI) & Collision steps (mean ± SD) & True minimum clearance (m) & Closed-loop ADE (m) & Closed-loop FDE (m) & Step time (ms) \\
\midrule
static & 12 & 83.3\% (55.2–95.3\%) & 0.50 ± 1.17 & 0.688 ± 0.579 & 2.103 ± 0.553 & 3.996 ± 1.103 & 95.12 ± 24.56 \\
cv & 12 & 75.0\% (46.8–91.1\%) & 2.42 ± 5.65 & 0.665 ± 0.923 & 0.490 ± 0.171 & 1.179 ± 0.433 & 96.17 ± 27.47 \\
kalman & 12 & 75.0\% (46.8–91.1\%) & 5.33 ± 9.95 & 0.528 ± 0.974 & 0.538 ± 0.188 & 1.251 ± 0.454 & 96.55 ± 25.92 \\
lstm & 12 & 100.0\% (75.8–100.0\%) & 0.00 ± 0.00 & 0.967 ± 0.788 & 0.570 ± 0.313 & 1.287 ± 0.707 & 99.00 ± 26.54 \\
\bottomrule
\end{tabularx}
\end{table}

The step-time means in Table 3 are calculated from the per-run MeanStepTime\_ms values. The complete P95 values and prediction/DWA/MPC/state-metric decomposition are retained in multiseed\_comparison.csv so that time steps are not treated as independent samples. LSTM succeeds in all 12 fixed-scene trials; static, cv and Kalman fail in 2, 3 and 3 trials, respectively. Failure types include collision|dwa\_no\_safe\_candidate, collision|goal\_not\_reached|dwa\_no\_safe\_candidate and goal\_not\_reached|dwa\_no\_safe\_candidate for cv in S3 seed 3. LSTM has no failure in this main fixed-scene suite.

LSTM has higher mean true minimum clearance than both traditional motion models and succeeds in all three S4 seeds. cv succeeds in only 1/3 S4 seeds and Kalman in 0/3. The fixed-window ADE of LSTM is not the lowest in every scene, whereas its closed-loop success and collision metrics are best. This indicates that closed-loop avoidance depends jointly on prediction, time alignment, the MPC safety cost, the flight-band constraint and DWA candidate screening rather than on a single ADE/FDE ranking.

\begin{table}[t]
\centering
\small
\caption{Timing decomposition for the fixed-scene main experiment (per-run mean re-aggregation, ms)}
\label{tab:table-4}
\begin{tabularx}{\linewidth}{YYYYYY}
\toprule
Method & Prediction & DWA & MPC & State/metrics & Step P95 \\
\midrule
static & 0.112 & 73.856 & 20.901 & 0.045 & 117.03 \\
cv & 0.122 & 73.654 & 22.132 & 0.045 & 111.39 \\
kalman & 3.366 & 71.751 & 21.211 & 0.043 & 111.21 \\
lstm & 4.375 & 73.263 & 21.115 & 0.043 & 117.26 \\
\bottomrule
\end{tabularx}
\end{table}

\begin{table}[t]
\centering
\small
\caption{Failure-type counts in the fixed-scene main experiment (12 trials per method)}
\label{tab:table-5}
\begin{tabularx}{\linewidth}{YYYYY}
\toprule
Method & Success & Collision + DWA no-safe-candidate & Collision + goal not reached + DWA no-safe-candidate & Goal not reached + DWA no-safe-candidate \\
\midrule
static & 10 & 2 & 0 & 0 \\
cv & 9 & 0 & 2 & 1 \\
kalman & 9 & 3 & 0 & 0 \\
lstm & 12 & 0 & 0 & 0 \\
\bottomrule
\end{tabularx}
\end{table}

\begin{table}[t]
\centering
\small
\caption{Scene-level fixed-scene closed-loop summary (successes/3; mean values)}
\label{tab:table-6}
\begin{tabularx}{\linewidth}{YYYYY}
\toprule
Scene & Method & Success/3 & Mean collision steps & Mean minimum clearance (m) \\
\midrule
S1 & static / cv / kalman / lstm & 3 / 3 / 3 / 3 & 0.00 / 0.00 / 0.00 / 0.00 & 1.209 / 1.938 / 1.930 / 2.136 \\
S2 & static / cv / kalman / lstm & 3 / 3 / 3 / 3 & 0.00 / 0.00 / 0.00 / 0.00 & 1.251 / 0.628 / 0.611 / 1.052 \\
S3 & static / cv / kalman / lstm & 3 / 2 / 3 / 3 & 0.00 / 0.00 / 0.00 / 0.00 & 0.298 / 0.387 / 0.231 / 0.262 \\
S4 & static / cv / kalman / lstm & 1 / 1 / 0 / 3 & 2.00 / 9.67 / 21.33 / 0.00 & −0.008 / −0.292 / −0.660 / 0.419 \\
\bottomrule
\end{tabularx}
\end{table}

\begin{figure}[t]
\centering
\includegraphics[width=0.98\linewidth]{../../results/figures_unified/S1 双障碍轻度交叉.png}
\caption{Three-dimensional closed-loop trajectories in S1}
\label{fig:figure-1}
\end{figure}

\begin{figure}[t]
\centering
\includegraphics[width=0.98\linewidth]{../../results/figures_unified/S2 四障碍走廊交错.png}
\caption{Three-dimensional closed-loop trajectories in S2}
\label{fig:figure-2}
\end{figure}

\begin{figure}[t]
\centering
\includegraphics[width=0.98\linewidth]{../../results/figures_unified/S3 五障碍非线性走廊压缩.png}
\caption{Three-dimensional closed-loop trajectories in S3}
\label{fig:figure-3}
\end{figure}

\begin{figure}[t]
\centering
\includegraphics[width=0.98\linewidth]{../../results/figures_unified/S4 六障碍密集走廊压缩.png}
\caption{Three-dimensional closed-loop trajectories in S4}
\label{fig:figure-4}
\end{figure}

\begin{figure}[t]
\centering
\includegraphics[width=0.98\linewidth]{../../results/figures_unified/S1 双障碍轻度交叉_pred.png}
\caption{Obstacle history and future-prediction comparison at a middle time in S1}
\label{fig:figure-5}
\end{figure}

\begin{figure}[t]
\centering
\includegraphics[width=0.98\linewidth]{../../results/figures_unified/S2 四障碍走廊交错_pred.png}
\caption{Obstacle history and future-prediction comparison at a middle time in S2}
\label{fig:figure-6}
\end{figure}

\begin{figure}[t]
\centering
\includegraphics[width=0.98\linewidth]{../../results/figures_unified/S4 六障碍密集走廊压缩_pred.png}
\caption{Obstacle history and future-prediction comparison at a middle time in S4}
\label{fig:figure-7}
\end{figure}

\begin{figure}[t]
\centering
\includegraphics[width=0.98\linewidth]{../../results/figures_unified/lstm_detour_evidence_S1.png}
\caption{LSTM detour and true-clearance evidence in S1}
\label{fig:figure-8}
\end{figure}

\begin{figure}[t]
\centering
\includegraphics[width=0.98\linewidth]{../../results/figures_unified/lstm_detour_evidence_S2.png}
\caption{LSTM detour and true-clearance evidence in S2}
\label{fig:figure-9}
\end{figure}

\begin{figure}[t]
\centering
\includegraphics[width=0.98\linewidth]{../../results/figures_unified/lstm_detour_evidence_S3.png}
\caption{LSTM detour and true-clearance evidence in S3}
\label{fig:figure-10}
\end{figure}

\begin{figure}[t]
\centering
\includegraphics[width=0.98\linewidth]{../../results/figures_unified/lstm_detour_evidence_S4.png}
\caption{LSTM detour and true-clearance evidence in S4}
\label{fig:figure-11}
\end{figure}

To complete the fixed-scene evidence, Figs. 8–11 report the LSTM three-dimensional true trajectory, along-track lateral offset, lateral/vertical offsets and true clearance for S1–S4, respectively. The \texttt{detours} count follows the code-defined event threshold of a ±0.5 m lateral offset; therefore, an event marker is not expected in every scene, and the S3 marker does not mean that evidence is missing for the other scenes. The random-scene and sensing-pressure suites are multi-run experiments; their primary evidence remains the run-level CSV files, summary tables, failure types and pressure-condition statistics rather than a single representative trajectory panel.

\subsection{Structural Ablation and Repeated Encounters}

In a single-seed diagnostic encounter, two-layer LSTM injection into both MPC and DWA with w\_route = 3 achieves a true minimum clearance of 1.939 m. Injection into MPC only and DWA only achieves 1.111 m and 1.921 m, respectively. The LSTM run with w\_route = 0 also succeeds in the same encounter, while the route-recovery cost is used to suppress persistent lateral offset after detours. Increasing the injection horizon from 10 to 20 and 30 steps changes the true minimum clearance from 0.858 to 1.727 and 1.939 m. Because this ablation uses one seed, these values diagnose mechanism rather than estimate stable overall effects.

\subsection{Observation-Noise and Delay Pressure Tests}

S4 is evaluated under clean observation, 0.05 m position noise, one-step delay and joint noise-plus-delay conditions. Each condition contains three matched seeds, and all methods share the same observation random stream. This pressure suite is one fixed six-obstacle S4 scene and 48 run-level records in total; it is not a four-scene or random-scene robustness estimate.

\begin{table}[t]
\centering
\small
\caption{S4 sensing-pressure results (three trials per method; means)}
\label{tab:table-7}
\begin{tabularx}{\linewidth}{YYYYYYY}
\toprule
Condition & Method & Success/3 & Collision steps & Minimum clearance (m) & ADE/FDE (m) & Step time (ms) \\
\midrule
clean & static & 1 & 2.33 & −0.006 & 2.575/4.839 & 126.62 \\
clean & cv & 0 & 14.00 & −0.779 & 0.764/1.884 & 161.90 \\
clean & kalman & 0 & 22.00 & −0.659 & 0.838/1.991 & 153.61 \\
clean & lstm & 2 & 0.00 & 0.365 & 0.684/1.554 & 166.13 \\
noise 0.05 m & static & 0 & 18.00 & −0.681 & 2.578/4.839 & 205.79 \\
noise 0.05 m & cv & 1 & 7.33 & −0.093 & 1.089/2.353 & 192.17 \\
noise 0.05 m & kalman & 2 & 5.33 & 0.202 & 0.956/2.152 & 197.84 \\
noise 0.05 m & lstm & 2 & 0.00 & 0.316 & 0.741/1.622 & 188.27 \\
delay 1 step & static & 0 & 6.00 & −0.120 & 2.734/4.980 & 132.48 \\
delay 1 step & cv & 0 & 23.00 & −0.665 & 0.828/1.979 & 124.15 \\
delay 1 step & kalman & 0 & 24.00 & −0.723 & 0.909/2.093 & 128.21 \\
delay 1 step & lstm & 2 & 0.00 & 0.770 & 0.731/1.625 & 138.40 \\
noise + delay & static & 0 & 6.67 & −0.129 & 2.738/4.981 & 134.61 \\
noise + delay & cv & 3 & 0.00 & 0.519 & 1.170/2.462 & 126.58 \\
noise + delay & kalman & 3 & 0.00 & 0.526 & 1.028/2.249 & 130.84 \\
noise + delay & lstm & 1 & 0.00 & 0.572 & 0.788/1.681 & 141.61 \\
\bottomrule
\end{tabularx}
\end{table}

\begin{table}[t]
\centering
\small
\caption{Timing decomposition for the S4 sensing-pressure tests (means, ms)}
\label{tab:table-8}
\begin{tabularx}{\linewidth}{YYYYYYY}
\toprule
Condition & Method & Prediction & DWA & MPC & State/metrics & Step P95 \\
\midrule
clean & static & 0.151 & 100.502 & 25.616 & 0.063 & 150.74 \\
clean & cv & 0.196 & 112.091 & 49.192 & 0.071 & 263.44 \\
clean & kalman & 5.329 & 104.989 & 42.926 & 0.066 & 186.59 \\
clean & lstm & 7.866 & 113.077 & 44.764 & 0.068 & 204.36 \\
noise 0.05 m & static & 0.696 & 132.632 & 71.904 & 0.099 & 267.18 \\
noise 0.05 m & cv & 0.635 & 122.617 & 68.375 & 0.091 & 244.70 \\
noise 0.05 m & kalman & 6.507 & 123.309 & 67.531 & 0.089 & 248.45 \\
noise 0.05 m & lstm & 9.422 & 121.902 & 56.461 & 0.082 & 260.57 \\
delay 1 step & static & 0.203 & 102.380 & 29.528 & 0.058 & 160.37 \\
delay 1 step & cv & 0.213 & 94.811 & 28.767 & 0.058 & 147.87 \\
delay 1 step & kalman & 4.735 & 94.583 & 28.593 & 0.057 & 148.85 \\
delay 1 step & lstm & 6.113 & 103.813 & 28.141 & 0.056 & 159.72 \\
noise + delay & static & 0.434 & 105.140 & 28.681 & 0.056 & 157.96 \\
noise + delay & cv & 0.435 & 97.794 & 28.000 & 0.056 & 153.98 \\
noise + delay & kalman & 4.922 & 98.420 & 27.210 & 0.056 & 154.01 \\
noise + delay & lstm & 6.481 & 105.600 & 29.202 & 0.057 & 163.35 \\
\bottomrule
\end{tabularx}
\end{table}

Under clean, noise-only and delay-only conditions, LSTM achieves 2/3 success with zero collisions. Under the joint condition, LSTM achieves 1/3, whereas cv and Kalman achieve 3/3. All LSTM pressure failures are goal\_not\_reached|dwa\_no\_safe\_candidate and contain no collision; traditional failures include both collision and no-safe-candidate outcomes. We therefore claim no universal superiority under joint sensing disturbances.

\subsection{Extended Random Scenes}

The extended suite uses scene seed 20260905 to generate 12 scenes. Each scene contains 3–8 obstacles and a mixture of straight, turn, sine-speed, acceleration, hover-then-go, zigzag and waypoint maneuvers. Each scene is evaluated with run seeds 1–3 for all four methods, producing 144 matched trials. The summary is checked field by field from the random-summary record.

\begin{table}[t]
\centering
\small
\caption{Closed-loop summary over the extended random scenes (means)}
\label{tab:table-9}
\begin{tabularx}{\linewidth}{YYYYYYYY}
\toprule
Method & Runs & Success rate (95\% Wilson CI) & Mean collision steps & Mean minimum clearance (m) & Mean ADE/FDE (m) & Mean step time (ms) & Failure-type count \\
\midrule
static & 36 & 72.2\% (26/36; 56.0–84.2\%) & 9.78 & 0.678 & 1.797/3.439 & 121.05 & collision + DWA no-safe-candidate: 10 \\
cv & 36 & 91.7\% (33/36; 78.2–97.1\%) & 2.00 & 0.749 & 0.397/1.026 & 121.27 & collision + DWA no-safe-candidate: 3 \\
kalman & 36 & 91.7\% (33/36; 78.2–97.1\%) & 2.03 & 0.733 & 0.443/1.098 & 125.04 & collision + DWA no-safe-candidate: 3 \\
lstm & 36 & 100.0\% (36/36; 90.4–100.0\%) & 0.00 & 1.108 & 0.327/0.764 & 127.33 & success: 36 \\
\bottomrule
\end{tabularx}
\end{table}

Across the 12 random scenes, LSTM has no failure, lower ADE/FDE and larger mean minimum clearance than cv and Kalman. Static has 10 collision|dwa\_no\_safe\_candidate failures, while cv and Kalman each have three. The Wilson intervals describe success-rate uncertainty and are not used as a significance test. The random suite supports transfer beyond S1–S4 within the fixed generator, but it does not replace real-sensor or hardware-in-the-loop validation.

\section{Discussion}

The evidence chain has three levels: fixed-window prediction error, fixed-scene closed-loop behavior and sensing-pressure boundaries. LSTM is slightly better than cv and Kalman on the aggregate fixed windows and is more advantageous in the nonlinear S3 corridor. In the fixed-scene closed-loop tests, LSTM succeeds in all 12 trials and all three S4 seeds without collision. The pressure tests show 2/3 success under noise-only and delay-only conditions but 1/3 under joint noise and delay, demonstrating that a lower prediction error does not imply a disturbance-independent closed-loop guarantee.

The engineering contribution is to use the predictor as a shared time-indexed interface between MPC and DWA. The route-recovery soft cost, safety-boundary penalty, prediction-based candidate rejection and two-layer injection form a closed-loop protection chain. The single-seed ablation is consistent with this mechanism, but it cannot attribute a stable quantitative gain to each component. The random-scene suite tests whether the effect is restricted to S1–S4; it does not replace hardware-in-the-loop or real-flight evaluation.

Four limitations remain. First, obstacle trajectories are generated synthetically and do not cover occlusion, missed detections, obstacle interaction or multi-agent games. Second, although the training and validation protocol separates trajectories and prevents normalization leakage, model-selection uncertainty is based on one training split and one training seed. Third, the clean main suite and the pressure suite show that joint noise and delay substantially reduce LSTM goal reachability, and cv/Kalman can outperform it under that condition. Fourth, the fixed-scene mean step time is approximately 95–99 ms, while pressure cases rise to approximately 126–206 ms; the present evidence does not support a stable 10 Hz real-time deployment claim. Further work should train with richer joint disturbances, missed detections and delay distributions and evaluate the safety constraints on a real or hardware-in-the-loop platform.

\section{Conclusions}

This paper presents a three-dimensional UAV dynamic-obstacle avoidance method that injects LSTM obstacle-trajectory predictions into both MPC and DWA. The evidence chain combines fixed-window prediction, a Kalman baseline, three-seed closed-loop comparisons, structural ablations, four sensing-pressure conditions and extended random scenes. The revised training protocol uses a trajectory-level split, training-only normalization statistics, correctly aligned history perturbation, early stopping and best-validation-loss selection. LSTM achieves 12/12 success in the fixed scenes with fixed-window ADE/FDE of 0.49/1.14 m. Under clean, noise-only, delay-only and joint conditions, its success rates are 2/3, 2/3, 2/3 and 1/3, respectively, and all pressure failures are collision-free. Across 12 random scenes and 36 matched trials, LSTM achieves 36/36 success with ADE/FDE of 0.327/0.764 m. The supported conclusion is that a correctly time-aligned LSTM forecast can improve MPC-DWA closed-loop safety and random-scene success under the tested simulation conditions, while joint sensing disturbances, model uncertainty, real-time deployment and real-flight validation remain open issues.

\section*{References}
\begin{thebibliography}{99}
\bibitem{ref1} Chang X C, Wang J Y, Dang S L, et al. Research on UAV path planning algorithm based on the MPC-DWA strategy [J/OL]. Electronics Optics \& Control, 1–9 [2026-09-07]. Available at: \url{https://link.cnki.net/urlid/41.1227.TN.20260325.1153.006}.

\bibitem{ref2} Teng F, Wang Y C, Yao Y H, Zhang K. Dynamic obstacle avoidance planning for UAVs based on deep reinforcement learning. Journal of Beijing University of Aeronautics and Astronautics, online first, 2025. \url{https://doi.org/10.13700/j.bh.1001-5965.2025.0084}.

\bibitem{ref3} Shi P L, et al. Dynamic obstacle avoidance control of intelligent vehicle on large-curvature roads considering trajectory prediction. Journal of Chang'an University (Natural Science Edition), 2024(4): 161–174. \url{https://doi.org/10.19721/j.cnki.1671-8879.2024.04.015}.

\bibitem{ref4} Hahn B. Enhancing obstacle avoidance in dynamic window approach via dynamic obstacle behavior prediction. Actuators. 2025;14(5):207. \url{https://doi.org/10.3390/act14050207}.

\bibitem{ref5} Yu H, Zhang Z. Dynamic obstacle avoidance algorithm for UUV based on physics-prior LSTM network model and dynamic window approach. Ocean Engineering. 2026;124503. \url{https://doi.org/10.1016/j.oceaneng.2026.124503}.

\bibitem{ref6} Fox D, Burgard W, Thrun S. The dynamic window approach to collision avoidance. IEEE Robotics \& Automation Magazine. 1997;4(1):23–33. \url{https://doi.org/10.1109/100.580977}.

\bibitem{ref7} Kalman R E. A new approach to linear filtering and prediction problems. Journal of Basic Engineering. 1960;82(1):35–45. \url{https://doi.org/10.1115/1.3662552}.

\bibitem{ref8} Hochreiter S, Schmidhuber J. Long short-term memory. Neural Computation. 1997;9(8):1735–1780. \url{https://doi.org/10.1162/neco.1997.9.8.1735}.

\bibitem{ref9} Alahi A, Goovaerts T, Ramanathan V, et al. Social LSTM: Human trajectory prediction in crowded spaces. In: Proceedings of the IEEE Conference on Computer Vision and Pattern Recognition. 2016:961–971. \url{https://doi.org/10.1109/CVPR.2016.110}.

\bibitem{ref10} Cao Z, Chi W. Dynamic obstacle avoidance of UAV using chance-constrained model predictive control. Optimal Control Applications and Methods. 2025;46(5):1914–1931. \url{https://doi.org/10.1002/oca.3298}.

\bibitem{ref11} Mayne D Q, Rawlings J B, Rao C V, Scokaert P O M. Constrained model predictive control: Stability and optimality. Automatica. 2000;36(6):789–814. \url{https://doi.org/10.1016/S0005-1098(99)00214-9}.

\bibitem{ref12} Hewing L, Wabersich K P, Menner M, Zeilinger M N. Learning-based model predictive control: Toward safe learning in control. Annual Review of Control, Robotics, and Autonomous Systems. 2020;3:269–296. \url{https://doi.org/10.1146/annurev-control-090419-075625}.

\bibitem{ref13} Andersson J A E, Gillis J, Horn G, Rawlings J B, Diehl M. CasADi: A software framework for nonlinear optimization and optimal control. Mathematical Programming Computation. 2019;11:1–36. \url{https://doi.org/10.1007/s12532-018-0139-4}.

\end{thebibliography}

\end{document}
